Consider the partial differential equation
\begin{equation}
\frac{\partial^2 u}{\partial x^2}
+
2\tan x
\frac{\partial^2 u}{\partial x\,\partial y}
-
\frac{\partial^2 u}{\partial y^2}
=
0
\label{90000037:pde}
\end{equation}
on the domain
\begin{equation}
\Omega
=
\{
(x,y)
\mid
0<x<1
\}.
\end{equation}
\begin{teilaufgaben}
\item
Is the differential operator in the equation \eqref{90000037:pde}
elliptic, parabolic or hyperbolic?
\item
Is the origin contained in the interval of the $y$-axis that is in the past
of the point $P=(\frac{\pi}4,0)$, if the future is in the direction
of the $x$-axis?
\end{teilaufgaben}

\begin{hinweis}
To answer the question in b), it is not necessary to solve a differential
equation.
A discussion of the sign of the slopes should be sufficient.
In case a solution needs to be computed,
the following integrals may turn out to be useful:
\begin{align*}
\int \tan x\,dx &= -\log\cos x + C
\\
\int \frac{dx}{\cos x} &= \frac12\log\frac{1+\sin x}{1-\sin x} + C
=
\operatorname{artanh}\sin x + C
\end{align*}
in the relevant $x$-intervals.
\end{hinweis}

\begin{loesung}
\begin{figure}
\centering
\includeagraphics[]{graph.pdf}
\caption{Characteristics through the point $P=(\frac{\pi}4,0)$ for
the partial differential equation \eqref{90000037:pde}.
\label{90000037:graph}}
\end{figure}
\begin{teilaufgaben}
\item
The symbol matrix of the differential operator of
\eqref{90000037:pde}
is
\begin{align*}
A
=
\begin{pmatrix}
   1   & \tan x \\
\tan x &  -1
\end{pmatrix}
\qquad\Rightarrow\qquad
\det A
&=
-1 - \tan^2 x
=
-(1+\tan^2x)
\\
&=
\displaystyle
-\frac{\sin^2x +\cos^2 x}{\cos^2x}
=
-
\frac{1}{\cos^2x}
<
0.
\end{align*}
The operator is thus hyperbolic.
\item
The differential equation for the characteristisc is
\[
\dot{y}^2 - 2\tan x\cdot \dot{x}\dot{y} - \dot{x}^2 = 0.
\]
Dividing by $\dot{x}$ gives the easier diffrential equation
\[
(y')^2 - 2\tan x\cdot y' - 1 = 0
\]
which can be solved using the quadratic formula
\begin{equation}
y'
=
\tan x \pm\!\sqrt{\tan^2 x+1}.
\label{90000037:dgl}
\end{equation}
Note that the right hand side does not depend on $y$, so the differential
equation can be solved by a simple integration.
The solutions are functions
\[
y_{\pm}(x)
=
\int
\tan x \pm\!\sqrt{\tan^2 x+1}
\,dx
+ C.
\]
Since for $x\in [0,1]$ we have
\[
\renewcommand{\arraycolsep}{2pt}
\begin{array}{rcccl}
y_+'(x) &=& \tan x +\!\sqrt{\tan^2 x+1} &>& 0
\\
y_-'(x) &=& \tan x -\!\sqrt{\tan^2 x+1} &<& 0,
\end{array}
\]
we conclude that the slope of the $y_+$ curve is positive,
while the slope of the $y_-$-curve is always negative.
Thus the $y_-$-curve through $P$ starts on the positive $y$-axis
and slopes down through the first quadrant,
while the $y_+$-curve through $P$ starts on the negative $y$-axis
and slopes up throught the fourth quadrant.
In particular, the origin is between those two characteristics
and it can in fact influence the value of the solution in $P$.
\qedhere
\end{teilaufgaben}
\end{loesung}

\begin{diskussion}
To decide whether the origin can influence the point $(\frac{\pi}4,0)$,
we did not need to determine the characteristics.
But using the integral formulas given in the hint, it is in fact possible
to give closed form expressions for the characteristics.

The square root on the right hand side of
\eqref{90000037:dgl}
can be written as
\begin{align*}
\!\sqrt{\tan^2 x+1}
&=
\!\sqrt{
\frac{\mathstrut\smash{\cos\!\mathstrut^2\,x+\sin\!\mathstrut^2\,x}}{\cos^2x}
}
=
\!\sqrt{\frac{1}{\cos^2x}}
=
\frac{1}{\cos x},
\end{align*}
since in $\Omega$ we have $\cos x > 0$.
Combining this with \eqref{90000037:dgl} gives the solution
\begin{align}
y(x)
&=
\int \tan x\pm\frac{1}{\cos x}\,dx
=
-\log\cos x \pm \log\frac{1 + \sin x}{1 - \sin x} + C.
\label{90000037:solution}
\end{align}

We can even to determine the value of $C$ for the characteristics
that pass through the point $(\frac{\pi}{4},0)$.
For this we need to substitute $x=\frac{\pi}{4}$ into the
right hand side of \eqref{90000037:solution}
and equate it to $0$.
Since
\[
\sin\frac{\pi}4=\cos\frac{\pi}4=\frac{\!\sqrt{2}}{2},
\]
this gives
\begin{align*}
0
&=
-\log\cos \frac{\pi}4
\pm
\log\frac{1+\sin\frac{\pi}4}{1-\sin\frac{\pi}4}
+
C_\pm
\\
C_\pm
&=
\log\frac{\!\sqrt{2}}{2}
\mp
\log\frac{1+\frac{\!\sqrt{2}}{2}}{1-\frac{\!\sqrt{2}}2}
=
-\log\!\sqrt{2}
\mp
\log\frac{2+\!\sqrt{2}}{2-\!\sqrt{2}}
=
-\log\!\sqrt{2}
\mp
\log\frac{\!\sqrt{2}+1}{\!\sqrt{2}-1}
\\
&=
\left\{
\begin{aligned}
-1.22794717729951567993\phantom{.}
\\
0.53479999673957037053.
\end{aligned}
\right.
\end{align*}
These are also the $y$-values at which the characteristics through
$P$ intersect the $y$-axis.
The past of the point $P$ is indicated in blue in 
figure~\ref{90000037:graph}.

Finally note that the product of the two slopes is by Vieta's theorem
\[
y'_+\cdot y'_- = -1,
\]
so the characteristics always intersect orthogonally.
This also is clearly visible in figure~\ref{90000037:graph}.
\end{diskussion}

\begin{bewertung}
Symbol matrix ({\bf A}) 1 point,
operator is hyperbolic ({\bf H}) 1 point,
differential equation of characteristics ({\bf D}) 1 point,
factorization or other method to convert them into two first order
differential equations ({\bf F}) 1 point,
characteristic curves or slopes ({\bf C}) 1 point,
is the point $P$ in the domain of influence ({\bf P}) 1 point.
\end{bewertung}

